\documentclass{article}
\usepackage[polish]{babel}
\usepackage[T1]{fontenc}
\usepackage[utf8]{inputenc}
\usepackage{amsmath}

\title{Wst\k{e}p do topologii algebraicznej}
\author{J.~Kowalski}
\date{2026}

\begin{document}
\maketitle

\section{Wst\k{e}p}

Topologia algebraiczna bada w\l{}a\'sno\'sci przestrzeni topologicznych za
pomoc\k{a} narz\k{e}dzi algebraicznych.  G\l{}\'ownym obiektem badania s\k{a}
grupy homotopii i grupy homologii.

Niech $X$ b\k{e}dzie przestrzeni\k{a} topologiczn\k{a} i niech $x_0 \in X$
b\k{e}dzie punktem bazowym.  Grup\k{e} podstawow\k{a} definiujemy jako
\[
  \pi_1(X, x_0) = \{[\gamma] : \gamma \text{ jest p\k{e}tl\k{a} w } x_0\}.
\]

\section{Twierdzenie g\l{}\'owne}

\begin{theorem}[Van Kampen]
Je\'sli $X = U_1 \cup U_2$, gdzie $U_1, U_2$ i $U_1 \cap U_2$ s\k{a}
\l{}\k{a}kowo sp\'ojne i otwarte, to $\pi_1(X)$ jest amalgamowanym produktem
wolnym $\pi_1(U_1) *_{\pi_1(U_1 \cap U_2)} \pi_1(U_2)$.
\end{theorem}

Dzi\k{e}ki temu twierdzeniu mo\.zna obliczy\'c grupy podstawowe wielu
przestrzeni, na przyk\l{}ad tor\'ow, sfer i powierzchni Riemanna.

\section{Podsumowanie}

Metody algebraiczne dostarczaj\k{a} pot\k{e}\.znych narz\k{e}dzi do badania
struktury przestrzeni topologicznych.

\end{document}
